\chapter{Los tres modelos}
\label{cap:modelos}

% GUION. Fuentes: TFM-Quantum/src/features.py, src/particion.py, doc/plan_eda.md,
%                 configs/gem_config.yaml
% ATENCION: el esqueleto anterior describia un modulo REM por GNN que esta FUERA DE ALCANCE.

\section{El problema de la representación}
\label{sec:representacion_modelos}
% ESTA seccion es el nucleo conceptual de la comparativa.
Cada muestra es una \textbf{matriz de altura variable} —de una decena a varios miles de filas,
una por puerta—. Un \textit{Graph Transformer} acepta esa entrada tal cual. Ridge y Random
Forest \textbf{no}: exigen que todas las filas tengan las mismas columnas.

\subsection{La agregación, y qué se pierde con ella}
Resumir el circuito en una fila \textbf{tira el orden y la estructura}: dos circuitos con las
mismas puertas en distinto orden dan \textbf{la misma fila}, y son físicamente distintos.

\textbf{Ahí está el núcleo de este trabajo:} la comparativa mide exactamente cuánto vale esa
información que el grafo conserva y la tabla pierde.

\section{Las variables agregadas}
\label{sec:features}
% Fuente: src/features.py (modulo de produccion, catalogo cerrado)
Ocho familias. Describir cada una en una frase y detallar solo las que resultaron predictivas:
\begin{enumerate}
    \item tamaño y forma del circuito;
    \item composición en tipos de puerta;
    \item ángulos de rotación;
    \item calidad de las puertas ejecutadas;
    \item telemetría del hardware que le tocó;
    \item cocientes físicos: duración frente a decoherencia;
    \item la heurística del estado del arte (el score de mapomatic como una variable más);
    \item tiempo.
\end{enumerate}

\subsection{Dos decisiones de diseño que importan}
% Fuente: src/features.py, docstrings de _telemetria_por_uso y _profundidad_y_camino
\textbf{Telemetría ponderada por uso:} un qubit malo que interviene en 200 puertas debe pesar
más que uno bueno usado dos veces.

\textbf{Duración crítica, no suma de duraciones:} las puertas sobre qubits distintos se
ejecutan simultáneamente. Lo que vive el estado es el \textbf{camino crítico del DAG}.

\subsection{Qué se excluye a propósito}
El \textbf{tipo de circuito} no es una variable: en entrenamiento es constante, y en test le
daría al modelo la respuesta al eje que precisamente se quiere medir.

\section{Ridge}
\label{sec:ridge}
\subsection{Por qué Ridge y no mínimos cuadrados}
% Fuente: doc/plan_eda.md seccion 4. Cifras del v1: REHACER.
Medido: la colinealidad entre variables es \textbf{extrema} —un bloque entero es «el tamaño»
escrito de ocho maneras—. Con columnas casi idénticas, mínimos cuadrados no puede decidir a
cuál atribuir el mérito: produce coeficientes enormes que \textbf{solo se cancelan dentro del
rango visto}. Y nuestro test es OOD en tres ejes.

\section{Random Forest}
\label{sec:rf}
El mejor baseline tabular según Liao et al. (2024). \textbf{Y sirve para algo más:} es inmune
a la colinealidad —cada árbol elige una columna y sigue—, así que la diferencia con Ridge no
es un problema a resolver, sino \textbf{parte de lo que la comparativa mide}.

\textbf{Su limitación, que hay que declarar:} \textbf{no puede extrapolar}. Satura en el valor
de la hoja más extrema que aprendió, y este conjunto de datos activa ese caso a propósito.

\section{El GEM: Graph Transformer}
\label{sec:gem}
\subsection{Por qué atención y no paso de mensajes}
El entrelazamiento crea dependencias entre puertas \textbf{muy separadas} en el grafo. Una GNN
convolucional necesitaría tantas capas como distancia; la atención las relaciona en una sola.

\subsection{El nodo virtual QCR}
Un nodo ficticio conectado a todos los demás: su embedding final resume el circuito entero y
produce la predicción, en lugar de un promedio sobre nodos. Respaldado por GTraQEM.

\subsection{La decisión sobre los grafos grandes}
% Fuente: doc/memoria/pendientes.md seccion 4
La atención es cuadrática y el peor caso no cabe en memoria. Presentar las tres opciones
—truncar, GPU grande, atención eficiente— y \textbf{justificar la elegida con la medida de a
cuántas muestras afecta}.

\section{El objetivo común}
\label{sec:objetivo_comun}
Los tres modelos predicen $f$ y reportan sobre $\Delta$ reconstruido, según el
capítulo~\ref{cap:objetivo}. \textbf{Es lo que hace que la comparación sea justa.}
